\documentclass[11pt]{article}
% ---------------------------------------------------------------
% R4.02 Qualite de developpement -- preambule commun des sujets de TP et des corrections
% Article clair, fond blanc, accents de couleur coherents avec les slides
% Charte reprise de R3.07 (meme palette IUT que les diapositives)
% Utilisation :
%   \documentclass[11pt]{article}
%   % ---------------------------------------------------------------
% R4.02 Qualite de developpement -- preambule commun des sujets de TP et des corrections
% Article clair, fond blanc, accents de couleur coherents avec les slides
% Charte reprise de R3.07 (meme palette IUT que les diapositives)
% Utilisation :
%   \documentclass[11pt]{article}
%   % ---------------------------------------------------------------
% R4.02 Qualite de developpement -- preambule commun des sujets de TP et des corrections
% Article clair, fond blanc, accents de couleur coherents avec les slides
% Charte reprise de R3.07 (meme palette IUT que les diapositives)
% Utilisation :
%   \documentclass[11pt]{article}
%   \input{../../latex/common/r402-tp-preamble.tex}
%   \begin{document}
%   \rtpheader{1}{Titre du TP}
% ---------------------------------------------------------------
\usepackage[T1]{fontenc}
\usepackage[utf8]{inputenc}
\usepackage[french]{babel}
\usepackage[scale=0.92]{tgheros}
\usepackage{tgpagella}
\usepackage[a4paper, margin=2.2cm, headheight=14pt]{geometry}
\usepackage{graphicx}
\usepackage{xcolor}
\usepackage{listings}
\usepackage[skins,listings,breakable]{tcolorbox}
\usepackage{enumitem}
\usepackage{booktabs}
\usepackage{array}
\usepackage{multirow}
\usepackage{tikz}
\usetikzlibrary{shapes, shapes.geometric, arrows, arrows.meta, positioning, calc}
\usepackage[colorlinks=true, linkcolor=qaaccentdark, urlcolor=qaaccentdark]{hyperref}
\usepackage{fancyhdr}
\usepackage{lastpage}
\usepackage{titlesec}
\usepackage{etoolbox}
\usepackage{eurosym}

% ---------------------------------------------------------------
% MENTION D'ELABORATION (discrete : gris clair, petit corps)
% ---------------------------------------------------------------
\definecolor{qaaimention}{HTML}{A3AAB5}
\newcommand{\raimention}{{\scriptsize\color{qaaimention}Supports \'elabor\'es avec l'aide de Claude Code Opus 5.5}}

% ---------------------------------------------------------------
% PALETTE -- version imprimable des couleurs des diapositives
% ---------------------------------------------------------------
\definecolor{qaaccentdark}{HTML}{7A1680}   % violet du logo IUT, fonce pour rester lisible sur blanc
\definecolor{qaaccent2dark}{HTML}{9B4B94}  % violet clair du logo IUT (accent secondaire)
\definecolor{qapanellight}{HTML}{F4F6F8}   % fond de bloc tres clair
\definecolor{qalinelight}{HTML}{D8DEE4}    % filets
\definecolor{qatextdark}{HTML}{1A1D23}     % texte principal
\definecolor{qamuteddark}{HTML}{5B6472}    % texte secondaire
\definecolor{qacodekwl}{HTML}{7A1680}      % code : mots-cles
\definecolor{qacodestrl}{HTML}{2E7D32}     % code : chaines
\definecolor{qacodecoml}{HTML}{6B7280}     % code : commentaires
\definecolor{qacodeannl}{HTML}{B26A00}     % code : annotations / javadoc

\color{qatextdark}

\setlength{\parskip}{4pt}
\setlength{\parindent}{0pt}

% ---------------------------------------------------------------
% TITRES DE SECTION -- texte blanc sur bande de couleur pleine largeur
% (gere aussi les variantes etoilees \section*)
% ---------------------------------------------------------------
\newcommand{\qasectionband}[2]{%
  \noindent\colorbox{#1}{%
    \parbox[c]{\dimexpr\linewidth-2\fboxsep\relax}{%
      \color{white}#2%
    }%
  }%
}
\newcommand{\secformat}[1]{\qasectionband{qaaccentdark}{\bfseries\Large\thesection\quad #1}}
\newcommand{\secformatnn}[1]{\qasectionband{qaaccentdark}{\bfseries\Large #1}}
\newcommand{\subsecformat}[1]{\qasectionband{qaaccent2dark}{\bfseries\large\thesubsection\quad #1}}
\newcommand{\subsecformatnn}[1]{\qasectionband{qaaccent2dark}{\bfseries\large #1}}

\titleformat{\section}{\normalfont}{}{0pt}{\secformat}
\titleformat{name=\section,numberless}{\normalfont}{}{0pt}{\secformatnn}
\titleformat{\subsection}{\normalfont}{}{0pt}{\subsecformat}
\titleformat{name=\subsection,numberless}{\normalfont}{}{0pt}{\subsecformatnn}
\titleformat{\subsubsection}{\normalfont\bfseries\color{qaaccentdark}}{\thesubsubsection}{0.6em}{}
\titlespacing*{\section}{0pt}{14pt}{8pt}
\titlespacing*{\subsection}{0pt}{10pt}{6pt}
\titlespacing*{\subsubsection}{0pt}{8pt}{4pt}

% ---------------------------------------------------------------
% LISTINGS -- theme clair (lstlisting reste utilisable tel quel)
% \lstcode : police + couleur de base, a reutiliser dans les surcharges
% ---------------------------------------------------------------
\newcommand{\lstcode}{\ttfamily\color{qatextdark}}

\lstdefinelanguage{TypeScript}{
  keywords={typeof, new, true, false, catch, function, return, null, switch, var,
    if, in, while, do, else, case, break, const, let, interface, class, extends,
    implements, enum, export, import, this, super, constructor, public, private,
    protected, static, async, await, throw, throws, try, finally, void},
  ndkeywords={boolean, any, number, string, int, long, double},
  sensitive=true,
  comment=[l]{//},
  morecomment=[s]{/*}{*/},
  morestring=[b]',
  morestring=[b]"
}

\lstdefinestyle{r402code}{
  basicstyle=\lstcode\footnotesize,
  keywordstyle=\color{qacodekwl}\bfseries,
  ndkeywordstyle=\color{qaaccent2dark}\bfseries,
  stringstyle=\color{qacodestrl},
  commentstyle=\color{qacodecoml}\itshape,
  morecomment=[s][\color{qacodeannl}]{/**}{*/},
  backgroundcolor=\color{qapanellight},
  rulecolor=\color{qalinelight},
  frame=single,
  framesep=4pt,
  xleftmargin=6pt,
  xrightmargin=6pt,
  numbers=left,
  numberstyle=\tiny\color{qamuteddark},
  numbersep=8pt,
  keepspaces=true,
  columns=fullflexible,
  showstringspaces=false,
  showspaces=false,
  showtabs=false,
  breaklines=true,
  breakatwhitespace=false,
  captionpos=b,
  tabsize=2,
  escapeinside={(*@}{@*)},
  % accents UTF-8 dans le code (listings ne gere pas l'UTF-8 multi-octets)
  literate={é}{{\'e}}1 {è}{{\`e}}1 {ê}{{\^e}}1 {à}{{\`a}}1 {â}{{\^a}}1
           {ù}{{\`u}}1 {û}{{\^u}}1 {ô}{{\^o}}1 {î}{{\^i}}1 {ç}{{\c c}}1
           {É}{{\'E}}1 {À}{{\`A}}1 {€}{{\euro{}}}1 {…}{{\ldots}}1,
}
\lstset{style=r402code}

% ---------------------------------------------------------------
% BOITES : exercice, code, information
% ---------------------------------------------------------------
\tcbset{
  exobox/.style={colback=qapanellight, colframe=qaaccentdark, boxrule=0.6pt,
    arc=2pt, left=8pt, right=8pt, top=6pt, bottom=6pt,
    coltitle=white, colbacktitle=qaaccentdark, fonttitle=\bfseries,
    boxsep=2pt, breakable},
  infobox/.style={colback=white, colframe=qaaccent2dark, boxrule=0.6pt,
    leftrule=3pt, arc=0pt, left=8pt, right=8pt, top=5pt, bottom=5pt,
    coltitle=qaaccent2dark, colbacktitle=white, fonttitle=\bfseries,
    breakable},
  codebox/.style={colback=qatextdark, colframe=qalinelight, boxrule=0.4pt,
    arc=2pt, left=8pt, right=8pt, top=5pt, bottom=5pt, listing only,
    listing options={basicstyle=\ttfamily\small\color{white},
      breaklines=true, showstringspaces=false, columns=fullflexible,
      keepspaces=true, frame=none, numbers=none, backgroundcolor={}}, breakable}
}
\newtcolorbox[auto counter]{exobox}[2][]{exobox,
  title={Exercice~\thetcbcounter\ifstrempty{#2}{}{ --- #2}}, #1}
\newtcolorbox{infobox}[1][\`A retenir]{infobox, title={#1}}
\newtcblisting{codebox}{codebox}

% ---------------------------------------------------------------
% EN-TETE DE PAGE (fil d'Ariane + pagination)
% ---------------------------------------------------------------
\newcommand{\qaheadtitlevalue}{}
\newcommand{\qaheadtitle}[1]{\renewcommand{\qaheadtitlevalue}{#1}}

\pagestyle{fancy}
\fancyhf{}
\renewcommand{\headrulewidth}{0.4pt}
\fancyhead[L]{\small\color{qamuteddark}R4.02 --- \qaheadtitlevalue}
\fancyhead[R]{\small\color{qamuteddark}\thepage\,/\,\pageref{LastPage}}
\fancyfoot[C]{\footnotesize\color{qamuteddark}IUT de Vannes --- Universit\'e Bretagne Sud}
\fancypagestyle{plain}{\fancyhf{}\renewcommand{\headrulewidth}{0pt}%
  \fancyfoot[C]{\footnotesize\color{qamuteddark}IUT de Vannes --- Universit\'e Bretagne Sud}}

% ---------------------------------------------------------------
% BANDEAU DE TITRE (pas une page dediee, pour rester compact)
% \rtpheader{numero ou intitule court}{titre du TP}
% ---------------------------------------------------------------
\newcommand{\ryear}{2026/2027}
\newcommand{\rtpheader}[2]{%
  \thispagestyle{plain}%
  \begin{minipage}[t]{0.68\textwidth}
    {\small\color{qamuteddark}
      Universit\'e Bretagne Sud\\
      IUT de Vannes\\
      D\'epartement Informatique --- \ryear}
  \end{minipage}%
  \hfill
  \begin{minipage}[t]{0.28\textwidth}
    \raggedleft
    \includegraphics[width=2.6cm]{iut_logo.jpg}
  \end{minipage}

  \vspace{6mm}
  \begin{center}
    {\large\bfseries\color{qatextdark} BUT 2 --- R4.02. Qualit\'e de d\'eveloppement}\\[2mm]
    {\normalsize\color{qamuteddark} TP #1}\\[1mm]
    {\Large\bfseries\color{qaaccentdark} #2}\\[2mm]
    \raimention
  \end{center}
  \vspace{4mm}
  \qaheadtitle{#2}
}

%   \begin{document}
%   \rtpheader{1}{Titre du TP}
% ---------------------------------------------------------------
\usepackage[T1]{fontenc}
\usepackage[utf8]{inputenc}
\usepackage[french]{babel}
\usepackage[scale=0.92]{tgheros}
\usepackage{tgpagella}
\usepackage[a4paper, margin=2.2cm, headheight=14pt]{geometry}
\usepackage{graphicx}
\usepackage{xcolor}
\usepackage{listings}
\usepackage[skins,listings,breakable]{tcolorbox}
\usepackage{enumitem}
\usepackage{booktabs}
\usepackage{array}
\usepackage{multirow}
\usepackage{tikz}
\usetikzlibrary{shapes, shapes.geometric, arrows, arrows.meta, positioning, calc}
\usepackage[colorlinks=true, linkcolor=qaaccentdark, urlcolor=qaaccentdark]{hyperref}
\usepackage{fancyhdr}
\usepackage{lastpage}
\usepackage{titlesec}
\usepackage{etoolbox}
\usepackage{eurosym}

% ---------------------------------------------------------------
% MENTION D'ELABORATION (discrete : gris clair, petit corps)
% ---------------------------------------------------------------
\definecolor{qaaimention}{HTML}{A3AAB5}
\newcommand{\raimention}{{\scriptsize\color{qaaimention}Supports \'elabor\'es avec l'aide de Claude Code Opus 5.5}}

% ---------------------------------------------------------------
% PALETTE -- version imprimable des couleurs des diapositives
% ---------------------------------------------------------------
\definecolor{qaaccentdark}{HTML}{7A1680}   % violet du logo IUT, fonce pour rester lisible sur blanc
\definecolor{qaaccent2dark}{HTML}{9B4B94}  % violet clair du logo IUT (accent secondaire)
\definecolor{qapanellight}{HTML}{F4F6F8}   % fond de bloc tres clair
\definecolor{qalinelight}{HTML}{D8DEE4}    % filets
\definecolor{qatextdark}{HTML}{1A1D23}     % texte principal
\definecolor{qamuteddark}{HTML}{5B6472}    % texte secondaire
\definecolor{qacodekwl}{HTML}{7A1680}      % code : mots-cles
\definecolor{qacodestrl}{HTML}{2E7D32}     % code : chaines
\definecolor{qacodecoml}{HTML}{6B7280}     % code : commentaires
\definecolor{qacodeannl}{HTML}{B26A00}     % code : annotations / javadoc

\color{qatextdark}

\setlength{\parskip}{4pt}
\setlength{\parindent}{0pt}

% ---------------------------------------------------------------
% TITRES DE SECTION -- texte blanc sur bande de couleur pleine largeur
% (gere aussi les variantes etoilees \section*)
% ---------------------------------------------------------------
\newcommand{\qasectionband}[2]{%
  \noindent\colorbox{#1}{%
    \parbox[c]{\dimexpr\linewidth-2\fboxsep\relax}{%
      \color{white}#2%
    }%
  }%
}
\newcommand{\secformat}[1]{\qasectionband{qaaccentdark}{\bfseries\Large\thesection\quad #1}}
\newcommand{\secformatnn}[1]{\qasectionband{qaaccentdark}{\bfseries\Large #1}}
\newcommand{\subsecformat}[1]{\qasectionband{qaaccent2dark}{\bfseries\large\thesubsection\quad #1}}
\newcommand{\subsecformatnn}[1]{\qasectionband{qaaccent2dark}{\bfseries\large #1}}

\titleformat{\section}{\normalfont}{}{0pt}{\secformat}
\titleformat{name=\section,numberless}{\normalfont}{}{0pt}{\secformatnn}
\titleformat{\subsection}{\normalfont}{}{0pt}{\subsecformat}
\titleformat{name=\subsection,numberless}{\normalfont}{}{0pt}{\subsecformatnn}
\titleformat{\subsubsection}{\normalfont\bfseries\color{qaaccentdark}}{\thesubsubsection}{0.6em}{}
\titlespacing*{\section}{0pt}{14pt}{8pt}
\titlespacing*{\subsection}{0pt}{10pt}{6pt}
\titlespacing*{\subsubsection}{0pt}{8pt}{4pt}

% ---------------------------------------------------------------
% LISTINGS -- theme clair (lstlisting reste utilisable tel quel)
% \lstcode : police + couleur de base, a reutiliser dans les surcharges
% ---------------------------------------------------------------
\newcommand{\lstcode}{\ttfamily\color{qatextdark}}

\lstdefinelanguage{TypeScript}{
  keywords={typeof, new, true, false, catch, function, return, null, switch, var,
    if, in, while, do, else, case, break, const, let, interface, class, extends,
    implements, enum, export, import, this, super, constructor, public, private,
    protected, static, async, await, throw, throws, try, finally, void},
  ndkeywords={boolean, any, number, string, int, long, double},
  sensitive=true,
  comment=[l]{//},
  morecomment=[s]{/*}{*/},
  morestring=[b]',
  morestring=[b]"
}

\lstdefinestyle{r402code}{
  basicstyle=\lstcode\footnotesize,
  keywordstyle=\color{qacodekwl}\bfseries,
  ndkeywordstyle=\color{qaaccent2dark}\bfseries,
  stringstyle=\color{qacodestrl},
  commentstyle=\color{qacodecoml}\itshape,
  morecomment=[s][\color{qacodeannl}]{/**}{*/},
  backgroundcolor=\color{qapanellight},
  rulecolor=\color{qalinelight},
  frame=single,
  framesep=4pt,
  xleftmargin=6pt,
  xrightmargin=6pt,
  numbers=left,
  numberstyle=\tiny\color{qamuteddark},
  numbersep=8pt,
  keepspaces=true,
  columns=fullflexible,
  showstringspaces=false,
  showspaces=false,
  showtabs=false,
  breaklines=true,
  breakatwhitespace=false,
  captionpos=b,
  tabsize=2,
  escapeinside={(*@}{@*)},
  % accents UTF-8 dans le code (listings ne gere pas l'UTF-8 multi-octets)
  literate={é}{{\'e}}1 {è}{{\`e}}1 {ê}{{\^e}}1 {à}{{\`a}}1 {â}{{\^a}}1
           {ù}{{\`u}}1 {û}{{\^u}}1 {ô}{{\^o}}1 {î}{{\^i}}1 {ç}{{\c c}}1
           {É}{{\'E}}1 {À}{{\`A}}1 {€}{{\euro{}}}1 {…}{{\ldots}}1,
}
\lstset{style=r402code}

% ---------------------------------------------------------------
% BOITES : exercice, code, information
% ---------------------------------------------------------------
\tcbset{
  exobox/.style={colback=qapanellight, colframe=qaaccentdark, boxrule=0.6pt,
    arc=2pt, left=8pt, right=8pt, top=6pt, bottom=6pt,
    coltitle=white, colbacktitle=qaaccentdark, fonttitle=\bfseries,
    boxsep=2pt, breakable},
  infobox/.style={colback=white, colframe=qaaccent2dark, boxrule=0.6pt,
    leftrule=3pt, arc=0pt, left=8pt, right=8pt, top=5pt, bottom=5pt,
    coltitle=qaaccent2dark, colbacktitle=white, fonttitle=\bfseries,
    breakable},
  codebox/.style={colback=qatextdark, colframe=qalinelight, boxrule=0.4pt,
    arc=2pt, left=8pt, right=8pt, top=5pt, bottom=5pt, listing only,
    listing options={basicstyle=\ttfamily\small\color{white},
      breaklines=true, showstringspaces=false, columns=fullflexible,
      keepspaces=true, frame=none, numbers=none, backgroundcolor={}}, breakable}
}
\newtcolorbox[auto counter]{exobox}[2][]{exobox,
  title={Exercice~\thetcbcounter\ifstrempty{#2}{}{ --- #2}}, #1}
\newtcolorbox{infobox}[1][\`A retenir]{infobox, title={#1}}
\newtcblisting{codebox}{codebox}

% ---------------------------------------------------------------
% EN-TETE DE PAGE (fil d'Ariane + pagination)
% ---------------------------------------------------------------
\newcommand{\qaheadtitlevalue}{}
\newcommand{\qaheadtitle}[1]{\renewcommand{\qaheadtitlevalue}{#1}}

\pagestyle{fancy}
\fancyhf{}
\renewcommand{\headrulewidth}{0.4pt}
\fancyhead[L]{\small\color{qamuteddark}R4.02 --- \qaheadtitlevalue}
\fancyhead[R]{\small\color{qamuteddark}\thepage\,/\,\pageref{LastPage}}
\fancyfoot[C]{\footnotesize\color{qamuteddark}IUT de Vannes --- Universit\'e Bretagne Sud}
\fancypagestyle{plain}{\fancyhf{}\renewcommand{\headrulewidth}{0pt}%
  \fancyfoot[C]{\footnotesize\color{qamuteddark}IUT de Vannes --- Universit\'e Bretagne Sud}}

% ---------------------------------------------------------------
% BANDEAU DE TITRE (pas une page dediee, pour rester compact)
% \rtpheader{numero ou intitule court}{titre du TP}
% ---------------------------------------------------------------
\newcommand{\ryear}{2026/2027}
\newcommand{\rtpheader}[2]{%
  \thispagestyle{plain}%
  \begin{minipage}[t]{0.68\textwidth}
    {\small\color{qamuteddark}
      Universit\'e Bretagne Sud\\
      IUT de Vannes\\
      D\'epartement Informatique --- \ryear}
  \end{minipage}%
  \hfill
  \begin{minipage}[t]{0.28\textwidth}
    \raggedleft
    \includegraphics[width=2.6cm]{iut_logo.jpg}
  \end{minipage}

  \vspace{6mm}
  \begin{center}
    {\large\bfseries\color{qatextdark} BUT 2 --- R4.02. Qualit\'e de d\'eveloppement}\\[2mm]
    {\normalsize\color{qamuteddark} TP #1}\\[1mm]
    {\Large\bfseries\color{qaaccentdark} #2}\\[2mm]
    \raimention
  \end{center}
  \vspace{4mm}
  \qaheadtitle{#2}
}

%   \begin{document}
%   \rtpheader{1}{Titre du TP}
% ---------------------------------------------------------------
\usepackage[T1]{fontenc}
\usepackage[utf8]{inputenc}
\usepackage[french]{babel}
\usepackage[scale=0.92]{tgheros}
\usepackage{tgpagella}
\usepackage[a4paper, margin=2.2cm, headheight=14pt]{geometry}
\usepackage{graphicx}
\usepackage{xcolor}
\usepackage{listings}
\usepackage[skins,listings,breakable]{tcolorbox}
\usepackage{enumitem}
\usepackage{booktabs}
\usepackage{array}
\usepackage{multirow}
\usepackage{tikz}
\usetikzlibrary{shapes, shapes.geometric, arrows, arrows.meta, positioning, calc}
\usepackage[colorlinks=true, linkcolor=qaaccentdark, urlcolor=qaaccentdark]{hyperref}
\usepackage{fancyhdr}
\usepackage{lastpage}
\usepackage{titlesec}
\usepackage{etoolbox}
\usepackage{eurosym}

% ---------------------------------------------------------------
% MENTION D'ELABORATION (discrete : gris clair, petit corps)
% ---------------------------------------------------------------
\definecolor{qaaimention}{HTML}{A3AAB5}
\newcommand{\raimention}{{\scriptsize\color{qaaimention}Supports \'elabor\'es avec l'aide de Claude Code Opus 5.5}}

% ---------------------------------------------------------------
% PALETTE -- version imprimable des couleurs des diapositives
% ---------------------------------------------------------------
\definecolor{qaaccentdark}{HTML}{7A1680}   % violet du logo IUT, fonce pour rester lisible sur blanc
\definecolor{qaaccent2dark}{HTML}{9B4B94}  % violet clair du logo IUT (accent secondaire)
\definecolor{qapanellight}{HTML}{F4F6F8}   % fond de bloc tres clair
\definecolor{qalinelight}{HTML}{D8DEE4}    % filets
\definecolor{qatextdark}{HTML}{1A1D23}     % texte principal
\definecolor{qamuteddark}{HTML}{5B6472}    % texte secondaire
\definecolor{qacodekwl}{HTML}{7A1680}      % code : mots-cles
\definecolor{qacodestrl}{HTML}{2E7D32}     % code : chaines
\definecolor{qacodecoml}{HTML}{6B7280}     % code : commentaires
\definecolor{qacodeannl}{HTML}{B26A00}     % code : annotations / javadoc

\color{qatextdark}

\setlength{\parskip}{4pt}
\setlength{\parindent}{0pt}

% ---------------------------------------------------------------
% TITRES DE SECTION -- texte blanc sur bande de couleur pleine largeur
% (gere aussi les variantes etoilees \section*)
% ---------------------------------------------------------------
\newcommand{\qasectionband}[2]{%
  \noindent\colorbox{#1}{%
    \parbox[c]{\dimexpr\linewidth-2\fboxsep\relax}{%
      \color{white}#2%
    }%
  }%
}
\newcommand{\secformat}[1]{\qasectionband{qaaccentdark}{\bfseries\Large\thesection\quad #1}}
\newcommand{\secformatnn}[1]{\qasectionband{qaaccentdark}{\bfseries\Large #1}}
\newcommand{\subsecformat}[1]{\qasectionband{qaaccent2dark}{\bfseries\large\thesubsection\quad #1}}
\newcommand{\subsecformatnn}[1]{\qasectionband{qaaccent2dark}{\bfseries\large #1}}

\titleformat{\section}{\normalfont}{}{0pt}{\secformat}
\titleformat{name=\section,numberless}{\normalfont}{}{0pt}{\secformatnn}
\titleformat{\subsection}{\normalfont}{}{0pt}{\subsecformat}
\titleformat{name=\subsection,numberless}{\normalfont}{}{0pt}{\subsecformatnn}
\titleformat{\subsubsection}{\normalfont\bfseries\color{qaaccentdark}}{\thesubsubsection}{0.6em}{}
\titlespacing*{\section}{0pt}{14pt}{8pt}
\titlespacing*{\subsection}{0pt}{10pt}{6pt}
\titlespacing*{\subsubsection}{0pt}{8pt}{4pt}

% ---------------------------------------------------------------
% LISTINGS -- theme clair (lstlisting reste utilisable tel quel)
% \lstcode : police + couleur de base, a reutiliser dans les surcharges
% ---------------------------------------------------------------
\newcommand{\lstcode}{\ttfamily\color{qatextdark}}

\lstdefinelanguage{TypeScript}{
  keywords={typeof, new, true, false, catch, function, return, null, switch, var,
    if, in, while, do, else, case, break, const, let, interface, class, extends,
    implements, enum, export, import, this, super, constructor, public, private,
    protected, static, async, await, throw, throws, try, finally, void},
  ndkeywords={boolean, any, number, string, int, long, double},
  sensitive=true,
  comment=[l]{//},
  morecomment=[s]{/*}{*/},
  morestring=[b]',
  morestring=[b]"
}

\lstdefinestyle{r402code}{
  basicstyle=\lstcode\footnotesize,
  keywordstyle=\color{qacodekwl}\bfseries,
  ndkeywordstyle=\color{qaaccent2dark}\bfseries,
  stringstyle=\color{qacodestrl},
  commentstyle=\color{qacodecoml}\itshape,
  morecomment=[s][\color{qacodeannl}]{/**}{*/},
  backgroundcolor=\color{qapanellight},
  rulecolor=\color{qalinelight},
  frame=single,
  framesep=4pt,
  xleftmargin=6pt,
  xrightmargin=6pt,
  numbers=left,
  numberstyle=\tiny\color{qamuteddark},
  numbersep=8pt,
  keepspaces=true,
  columns=fullflexible,
  showstringspaces=false,
  showspaces=false,
  showtabs=false,
  breaklines=true,
  breakatwhitespace=false,
  captionpos=b,
  tabsize=2,
  escapeinside={(*@}{@*)},
  % accents UTF-8 dans le code (listings ne gere pas l'UTF-8 multi-octets)
  literate={é}{{\'e}}1 {è}{{\`e}}1 {ê}{{\^e}}1 {à}{{\`a}}1 {â}{{\^a}}1
           {ù}{{\`u}}1 {û}{{\^u}}1 {ô}{{\^o}}1 {î}{{\^i}}1 {ç}{{\c c}}1
           {É}{{\'E}}1 {À}{{\`A}}1 {€}{{\euro{}}}1 {…}{{\ldots}}1,
}
\lstset{style=r402code}

% ---------------------------------------------------------------
% BOITES : exercice, code, information
% ---------------------------------------------------------------
\tcbset{
  exobox/.style={colback=qapanellight, colframe=qaaccentdark, boxrule=0.6pt,
    arc=2pt, left=8pt, right=8pt, top=6pt, bottom=6pt,
    coltitle=white, colbacktitle=qaaccentdark, fonttitle=\bfseries,
    boxsep=2pt, breakable},
  infobox/.style={colback=white, colframe=qaaccent2dark, boxrule=0.6pt,
    leftrule=3pt, arc=0pt, left=8pt, right=8pt, top=5pt, bottom=5pt,
    coltitle=qaaccent2dark, colbacktitle=white, fonttitle=\bfseries,
    breakable},
  codebox/.style={colback=qatextdark, colframe=qalinelight, boxrule=0.4pt,
    arc=2pt, left=8pt, right=8pt, top=5pt, bottom=5pt, listing only,
    listing options={basicstyle=\ttfamily\small\color{white},
      breaklines=true, showstringspaces=false, columns=fullflexible,
      keepspaces=true, frame=none, numbers=none, backgroundcolor={}}, breakable}
}
\newtcolorbox[auto counter]{exobox}[2][]{exobox,
  title={Exercice~\thetcbcounter\ifstrempty{#2}{}{ --- #2}}, #1}
\newtcolorbox{infobox}[1][\`A retenir]{infobox, title={#1}}
\newtcblisting{codebox}{codebox}

% ---------------------------------------------------------------
% EN-TETE DE PAGE (fil d'Ariane + pagination)
% ---------------------------------------------------------------
\newcommand{\qaheadtitlevalue}{}
\newcommand{\qaheadtitle}[1]{\renewcommand{\qaheadtitlevalue}{#1}}

\pagestyle{fancy}
\fancyhf{}
\renewcommand{\headrulewidth}{0.4pt}
\fancyhead[L]{\small\color{qamuteddark}R4.02 --- \qaheadtitlevalue}
\fancyhead[R]{\small\color{qamuteddark}\thepage\,/\,\pageref{LastPage}}
\fancyfoot[C]{\footnotesize\color{qamuteddark}IUT de Vannes --- Universit\'e Bretagne Sud}
\fancypagestyle{plain}{\fancyhf{}\renewcommand{\headrulewidth}{0pt}%
  \fancyfoot[C]{\footnotesize\color{qamuteddark}IUT de Vannes --- Universit\'e Bretagne Sud}}

% ---------------------------------------------------------------
% BANDEAU DE TITRE (pas une page dediee, pour rester compact)
% \rtpheader{numero ou intitule court}{titre du TP}
% ---------------------------------------------------------------
\newcommand{\ryear}{2026/2027}
\newcommand{\rtpheader}[2]{%
  \thispagestyle{plain}%
  \begin{minipage}[t]{0.68\textwidth}
    {\small\color{qamuteddark}
      Universit\'e Bretagne Sud\\
      IUT de Vannes\\
      D\'epartement Informatique --- \ryear}
  \end{minipage}%
  \hfill
  \begin{minipage}[t]{0.28\textwidth}
    \raggedleft
    \includegraphics[width=2.6cm]{iut_logo.jpg}
  \end{minipage}

  \vspace{6mm}
  \begin{center}
    {\large\bfseries\color{qatextdark} BUT 2 --- R4.02. Qualit\'e de d\'eveloppement}\\[2mm]
    {\normalsize\color{qamuteddark} TP #1}\\[1mm]
    {\Large\bfseries\color{qaaccentdark} #2}\\[2mm]
    \raimention
  \end{center}
  \vspace{4mm}
  \qaheadtitle{#2}
}

\graphicspath{{../../latex/common/assets/}{./}}

% A ajuster avant diffusion
\newcommand{\rstarterurl}{\url{https://github.com/ctiber/r402}}
\newcommand{\rdatelimite}{la veille de la séance du TP2, à 23\,h\,59}

\begin{document}

\rtpheader{1}{Projet fil rouge et JUnit avancé}

\begin{infobox}[Objectifs]
\begin{itemize}[leftmargin=1.2em, itemsep=1pt]
  \item Prendre en main le projet fil rouge que vous utiliserez pendant les six TP.
  \item Choisir des cas de test à partir d'une spécification : partitions d'équivalence et valeurs limites.
  \item Écrire des tests JUnit lisibles : Arrange / Act / Assert, tests paramétrés, \texttt{@Nested}, \texttt{assertAll}, AssertJ.
  \item Trouver les bugs du code grâce aux tests, les corriger, et protéger le code contre les régressions.
\end{itemize}
\end{infobox}

\section*{Contexte}

Vous reprenez l'application bancaire de R3.07 (clients, comptes, opérations, virements), réécrite avec Spring Boot et un front Vue 3. Un projet de départ vous est fourni. Il fonctionne, mais \textbf{il est peu testé et contient des bugs}. Au fil des TP, vous allez le tester à tous les niveaux : domaine (TP1), acceptation (TP2), services (TP3), intégration avec la base (TP4), bout en bout (TP5), puis automatiser le tout dans une pipeline (TP6).

Aujourd'hui, on se concentre sur la classe du domaine la plus importante : \texttt{Compte}. Elle ne dépend ni de la base de données, ni de Spring : on peut la tester avec JUnit seul.

Vous travaillez \textbf{en binôme}, sur un dépôt Git par binôme, pendant tout le module.

% ================================================================
\section{Mise en place (15 min)}

\begin{exobox}{Récupérer et lancer le projet}
\begin{enumerate}[leftmargin=1.4em]
  \item Un seul membre du binôme importe le projet de départ dans un nouveau projet \textbf{privé} de la forge GitLab de l'IUT (\textit{New project} $\rightarrow$ \textit{Import project} $\rightarrow$ \textit{Repository by URL}), à partir de l'adresse :\\ \rstarterurl\\
  L'import reprend toutes les branches du dépôt, dont \texttt{evolution-frais}, utilisée à la fin du TP.
  \item Il ajoute son binôme et l'enseignant comme membres du projet (rôle \textit{Developer} pour le binôme, \textit{Reporter} pour l'enseignant).
  \item Chacun clone le dépôt et l'ouvre dans IntelliJ IDEA.
  \item Lancez les tests existants. Ils doivent tous passer :
\begin{codebox}
./mvnw test
\end{codebox}
  \item Démarrez la base de données, puis l'application, et ouvrez \url{http://localhost:8080} :
\begin{codebox}
docker compose up -d
./mvnw spring-boot:run
\end{codebox}
\end{enumerate}
\end{exobox}

\begin{exobox}{Découvrir l'architecture}
Parcourez le projet et répondez aux questions suivantes dans un fichier \texttt{docs/tp1.md} (quelques lignes suffisent).
\begin{enumerate}[leftmargin=1.4em]
  \item Dans quels paquetages se trouvent le domaine, les services, les repositories et les contrôleurs REST ?
  \item De quels autres objets la classe \texttt{VirementService} a-t-elle besoin pour fonctionner ? Lesquels seraient gênants dans un test unitaire (lents, non reproductibles, extérieurs à l'application) ?
  \item Combien de tests contient le projet ? Quelles classes testent-ils ? Que se passerait-il si on modifiait \texttt{Compte.debiter} demain ?
  \item Où sont déclarées, dans le \texttt{pom.xml}, les dépendances JUnit et AssertJ ? Quelle est leur portée (\texttt{scope}) et pourquoi ?
\end{enumerate}
\end{exobox}

% ================================================================
\section{La spécification de \texttt{Compte}}

Voici l'interface publique de la classe \texttt{Compte} (paquetage \texttt{fr.iutvannes.banque.domaine}) :

\begin{lstlisting}[language=Java]
public class Compte {
    public static final BigDecimal PLAFOND_EPARGNE = new BigDecimal("22950.00");

    public Compte(String numero, BigDecimal soldeInitial)                  // compte courant, en euros
    public Compte(String numero, TypeCompte type, BigDecimal soldeInitial) // COURANT ou EPARGNE
    public Compte(String numero, Devise devise, BigDecimal soldeInitial)   // (utilise a partir du TP3)

    public void setDecouvertAutorise(BigDecimal decouvert)
    public void crediter(BigDecimal montant)
    public void debiter(BigDecimal montant)
    public void cloturer()

    public String getNumero()
    public TypeCompte getType()
    public BigDecimal getSolde()
    public BigDecimal getDecouvertAutorise()
    public boolean isCloture()
    public List<Operation> getOperations()   // historique, non modifiable
}
\end{lstlisting}

Les règles métier à respecter sont les suivantes. \textbf{C'est la référence} : si le code ne fait pas ce qui est écrit ici, c'est le code qui a tort.

\begin{center}
\renewcommand{\arraystretch}{1.3}
\begin{tabular}{@{}>{\bfseries\color{qaaccentdark}}l p{14cm}@{}}
\toprule
R1 & \textbf{Montant.} Le montant d'un crédit ou d'un débit est strictement positif et a au plus deux décimales. Sinon (y compris s'il vaut \texttt{null}) : \texttt{IllegalArgumentException}. \\
R2 & \textbf{Débit et découvert.} Un débit est accepté si le solde après l'opération reste supérieur ou égal à l'opposé du découvert autorisé (\texttt{solde - montant >= -decouvertAutorise}). Sinon : \texttt{SoldeInsuffisantException}. Le découvert autorisé vaut 0 par défaut ; il ne peut pas être négatif (\texttt{IllegalArgumentException}) ; il ne peut être modifié que sur un compte courant (\texttt{IllegalStateException} pour un compte épargne). \\
R3 & \textbf{Plafond épargne.} Le solde d'un compte épargne ne peut pas dépasser \texttt{PLAFOND\_EPARGNE} (22\,950,00\,\euro{}). Un crédit qui ferait dépasser ce plafond est refusé : \texttt{PlafondDepasseException}. Un compte courant n'a pas de plafond. \\
R4 & \textbf{Clôture.} Un compte ne peut être clôturé que si son solde est exactement nul ; sinon : \texttt{IllegalStateException}. Sur un compte clôturé, tout crédit ou débit lève \texttt{CompteClotureException}. \\
R5 & \textbf{Historique.} Chaque crédit ou débit \textbf{réussi} ajoute une \texttt{Operation} (type, montant, date) à l'historique. Une opération \textbf{refusée} ne modifie ni le solde, ni l'historique. \\
\bottomrule
\end{tabular}
\end{center}

\begin{infobox}[Utilitaire fourni]
La classe \texttt{TestUtils} (dans \texttt{src/test/java}) fournit \texttt{eur("12.50")}, un raccourci pour \texttt{new BigDecimal("12.50")}. Utilisez-la avec \texttt{import static fr.iutvannes.banque.TestUtils.eur;}.

Pour comparer des \texttt{BigDecimal}, utilisez \texttt{assertThat(x).isEqualByComparingTo("70")} (AssertJ) : \texttt{assertEquals} échoue entre \texttt{70} et \texttt{70.00} (voir le cours).
\end{infobox}

% ================================================================
\section{Choisir ses cas de test (15 min)}

\begin{exobox}{Partitions d'équivalence et valeurs limites}
Avant d'écrire la moindre ligne de code, complétez dans \texttt{docs/tp1.md} un tableau des cas de test pour les règles \textbf{R1, R2 et R3}. Pour chaque règle, identifiez les \textbf{partitions} (classes d'entrées qui se comportent pareil) et les \textbf{valeurs limites} (de chaque côté de chaque frontière). Les premières lignes sont données :

\medskip
{\footnotesize
\begin{tabular}{@{}lllll@{}}
\toprule
\textbf{Règle} & \textbf{Partition / limite} & \textbf{Contexte} & \textbf{Entrée} & \textbf{Résultat attendu} \\
\midrule
R1 & montant nul (limite) & solde 100.00 & \texttt{debiter(0)} & \texttt{IllegalArgumentException} \\
R1 & plus petit montant valide (limite) & solde 100.00 & \texttt{debiter(0.01)} & solde 99.99 \\
R1 & trop de décimales & solde 100.00 & \texttt{crediter(10.001)} & \texttt{IllegalArgumentException} \\
R2 & \ldots & \ldots & \ldots & \ldots \\
\bottomrule
\end{tabular}}
\medskip

\textit{Indice} : pour R2, pensez au cas où le découvert autorisé est 0 \textbf{et} au cas où il vaut par exemple 200.00. Pour R3, où se trouve exactement la frontière ?
\end{exobox}

% ================================================================
\section{Écrire les tests (40 min)}

Écrivez vos tests dans \texttt{src/test/java/fr/iutvannes/banque/domaine/CompteTest.java}. Respectez la structure Arrange / Act / Assert et donnez à chaque test un nom (ou un \texttt{@DisplayName}) qui dit ce qu'il vérifie.

\begin{exobox}{Tests paramétrés pour R1 et R2}
\begin{enumerate}[leftmargin=1.4em]
  \item Transformez les lignes de votre tableau pour R1 en \textbf{un} test paramétré pour les montants invalides (\texttt{@ValueSource} ou \texttt{@CsvSource}) et un autre pour les montants valides. Le cas \texttt{null} se traite avec \texttt{@NullSource}.
  \item Faites de même pour R2, avec \texttt{@CsvSource} : solde initial, découvert autorisé, montant, résultat attendu. Faut-il un seul test paramétré, ou un pour les débits acceptés et un pour les débits refusés ? Justifiez en une phrase dans \texttt{docs/tp1.md}.
  \item Pour les débits refusés, vérifiez avec \texttt{assertThrows} (ou \texttt{assertThatThrownBy}) le type de l'exception \textbf{et} que le solde n'a pas changé.
  \item Utilisez l'attribut \texttt{name} de \texttt{@ParameterizedTest} pour que le rapport de tests soit lisible.
\end{enumerate}
\end{exobox}

\begin{exobox}{Organiser les tests par contexte avec \texttt{@Nested}}
Ajoutez dans \texttt{CompteTest} des classes imbriquées (\texttt{@Nested} + \texttt{@DisplayName}), chacune avec son propre \texttt{@BeforeEach} :
\begin{itemize}[leftmargin=1.2em]
  \item \textbf{« Un compte épargne »} : tests de R3 (y compris : un compte courant n'a pas de plafond ; le découvert d'un compte épargne ne peut pas être modifié) ;
  \item \textbf{« Un compte clôturé »} : tests de R4 (clôture d'un compte à solde nul, positif, \textbf{négatif} ; opérations sur un compte clôturé) ;
  \item \textbf{« L'historique »} : tests de R5. Après un crédit puis un débit réussis, vérifiez en un seul test, avec \texttt{assertAll} ou AssertJ (\texttt{extracting}, \texttt{containsExactly}), le nombre d'opérations, leurs types et leurs montants. Vérifiez aussi qu'un débit \textbf{refusé} n'ajoute rien.
\end{itemize}
Lancez les tests depuis IntelliJ et vérifiez que l'arborescence affichée se lit comme une spécification.
\end{exobox}

% ================================================================
\section{Corriger les bugs, sans régression (20 min)}

Normalement, plusieurs de vos tests sont \textbf{rouges}. Avant de toucher au code, assurez-vous que c'est bien le code qui a tort et pas le test : relisez la règle concernée.

\begin{exobox}{Corriger un bug à la fois}
Pour \textbf{chaque} bug trouvé :
\begin{enumerate}[leftmargin=1.4em]
  \item notez dans \texttt{docs/tp1.md} : la règle violée, le test qui le révèle, la cause dans le code ;
  \item corrigez le code de \texttt{Compte} (et seulement lui : on ne modifie pas un test pour le faire passer) ;
  \item relancez \textbf{tous} les tests, pas seulement celui-ci ;
  \item faites un commit dont le message cite le bug corrigé.
\end{enumerate}
Si aucun de vos tests n'échoue, c'est mauvais signe : complétez votre tableau de cas, il manque des valeurs limites.
\end{exobox}

\begin{exobox}{Non-régression : la branche d'un collègue}
Un collègue a travaillé en parallèle sur la branche \texttt{evolution-frais} : il a ajouté les frais de tenue de compte et en a profité pour « simplifier » le code de \texttt{Compte}. Il affirme que « tout marche, l'application démarre ».
\begin{enumerate}[leftmargin=1.4em]
  \item Fusionnez sa branche dans la vôtre :
\begin{codebox}
git fetch origin
git merge origin/evolution-frais
\end{codebox}
  \item Relancez tous les tests. Quelle règle n'est plus respectée ? Qu'est-ce qui l'aurait révélé si vos tests n'avaient pas existé ?
  \item Corrigez la régression. Expliquez en deux ou trois phrases dans \texttt{docs/tp1.md} en quoi votre suite de tests a servi de filet de sécurité.
  \item \textit{Pour aller plus loin} : la nouvelle méthode \texttt{appliquerFraisTenueCompte} n'est pas testée. Écrivez sa spécification (en vous inspirant des règles R1 à R5), puis ses tests.
\end{enumerate}
\end{exobox}

% ================================================================
\section{Rendu}

\begin{infobox}[À rendre avant \rdatelimite]
Sur le dépôt du binôme, un \textbf{tag} \texttt{tp1} pointant sur le commit rendu :
\begin{codebox}
git tag tp1
git push origin tp1
\end{codebox}
Le commit doit contenir :
\begin{itemize}[leftmargin=1.2em, itemsep=1pt]
  \item \texttt{CompteTest.java} : tests paramétrés (R1, R2), contextes \texttt{@Nested} (R3, R4, R5) ;
  \item \texttt{Compte.java} corrigé : \texttt{./mvnw test} passe entièrement ;
  \item \texttt{docs/tp1.md} : réponses sur l'architecture, tableau des cas de test, journal des bugs, analyse de la régression.
\end{itemize}
Un rendu manquant ou en retard compte comme un défaut pour le coefficient d'assiduité (voir le cours).
\end{infobox}

\end{document}
